\chapter{Budd--Chiari Syndrome}

\section*{Definition and Overview}
Budd--Chiari syndrome is a rare disorder in which the veins draining blood from the liver become blocked, usually by a blood clot, causing obstruction of hepatic venous outflow. The blockage raises pressure within the liver, leading to enlargement and congestion of the organ, abdominal pain, accumulation of fluid in the abdomen (ascites), and, if uncorrected, progressive liver damage and failure. The condition may present acutely, subacutely, or with features of chronic liver disease, and it requires specialist hepatology care. With modern treatment, including anticoagulation and interventional or surgical procedures, many affected people can be stabilised and live well for many years.

\section*{Epidemiology}
Budd--Chiari syndrome is uncommon, affecting roughly one to two people per million per year. It can occur at any age but is most often diagnosed in young adults, and women are affected somewhat more often than men. In parts of Asia, obstruction by a congenital fibrous web or membrane within the inferior vena cava is a more common cause than in Western populations. Most affected people have an underlying thrombotic tendency, frequently an overt or occult myeloproliferative disorder, which may not be recognised before the diagnosis.

\section*{Aetiology and Pathophysiology}
The syndrome results from obstruction of the hepatic veins or of the intrahepatic segment of the inferior vena cava. When outflow is blocked, blood cannot leave the liver normally, pressure rises in the hepatic sinusoids, and the liver becomes congested, enlarged, and tender. Prolonged congestion damages hepatocytes and eventually leads to fibrosis and cirrhosis.

Blockage is most often caused by a thrombus, and an underlying clotting tendency is found in the majority of cases.

\begin{itemize}
    \item \textbf{Myeloproliferative disorders:} polycythaemia vera and essential thrombocythaemia are the most commonly identified causes
    \item \textbf{Inherited thrombophilia:} factor V Leiden, prothrombin gene mutations, and protein C, protein S, or antithrombin deficiency
    \item \textbf{Acquired thrombophilia:} antiphospholipid syndrome and paroxysmal nocturnal haemoglobinuria
    \item \textbf{Hormonal factors:} pregnancy and the combined oral contraceptive pill increase thrombotic risk
    \item \textbf{Inflammatory disease:} Behcet's disease and other vasculitides
    \item \textbf{Local causes:} tumour invasion of the vein or a caval membrane or web
    \item \textbf{No cause identified:} in a minority of cases no underlying disorder is found
\end{itemize}

\section*{Clinical Features}
Symptoms may develop suddenly or evolve gradually over weeks to months, depending on the speed of obstruction and the development of collateral drainage.

\begin{itemize}
    \item Abdominal pain and discomfort, particularly in the right upper quadrant
    \item Abdominal swelling from ascites
    \item An enlarged, tender liver palpable on examination
    \item Nausea, vomiting, and loss of appetite
    \item Fatigue and generalised weakness
    \item Mild jaundice (yellowing of the skin and eyes)
    \item In severe disease, bleeding from oesophageal varices and hepatic encephalopathy with confusion
\end{itemize}

\section*{Differential Diagnosis}
Several other liver and abdominal conditions can present with similar features and must be excluded.

\begin{itemize}
    \item \textbf{Other causes of acute liver failure:} paracetamol overdose, viral hepatitis, and drug-induced liver injury
    \item \textbf{Cirrhosis from other causes:} alcohol-related liver disease and non-alcoholic fatty liver disease
    \item \textbf{Cardiac causes of hepatic congestion:} right heart failure and constrictive pericarditis
    \item \textbf{Biliary disease:} gallstones, cholecystitis, or obstruction of the bile ducts
    \item \textbf{Liver masses or infection:} tumour, hepatic abscess, or hydatid disease
    \item \textbf{Portal vein thrombosis:} obstruction of the main vein entering the liver, a distinct condition with similar features
    \item \textbf{Other causes of ascites:} malignancy, tuberculosis, or gynaecological disease
\end{itemize}

\section*{When to Seek Medical Care}
See a doctor promptly for persistent right upper abdominal pain with a swollen abdomen, particularly in the setting of a known clotting disorder, pregnancy, or combined oral contraceptive use. The diagnosis of Budd--Chiari syndrome is usually made in hospital.

\textbf{Call 000 (or local emergency services) for:}
\begin{itemize}
    \item Vomiting blood or passing black, tarry stools
    \item Confusion, drowsiness, or a change in behaviour, suggesting liver failure
    \item Sudden severe abdominal pain with a rapidly swelling abdomen
    \item Difficulty breathing from a very swollen abdomen
    \item Any sudden deterioration in someone with known liver disease
\end{itemize}

\section*{Diagnosis and Investigations}
Imaging of the hepatic veins is central to the diagnosis, supported by laboratory testing.

\begin{itemize}
    \item \textbf{Ultrasound with Doppler:} the first-line test, demonstrating the hepatic veins and detecting thrombus or absent flow
    \item \textbf{CT or MRI scan:} detailed cross-sectional imaging of the liver, its vasculature, and any tumour
    \item \textbf{Blood tests:} liver function tests, clotting studies, and a full blood count
    \item \textbf{Thrombophilia screen:} testing for inherited and acquired thrombotic disorders
    \item \textbf{Liver biopsy:} sometimes performed to assess the degree of fibrosis and congestion
    \item \textbf{Venography:} a dye study of the hepatic veins and inferior vena cava, often combined with therapeutic intervention
\end{itemize}

\section*{Management}
Treatment aims to restore hepatic venous outflow, prevent further thrombosis, and manage the complications of portal hypertension.

\begin{itemize}
    \item \textbf{Anticoagulation:} blood-thinning medicines to prevent extension of the clot and new thromboses
    \item \textbf{Treatment of the underlying cause:} managing an identified myeloproliferative or thrombophilic disorder and withdrawing medicines that raise clot risk
    \item \textbf{Interventional radiology:} angioplasty or stenting to reopen a narrowed or blocked vein, or insertion of a transjugular intrahepatic portosystemic shunt (TIPSS) to reroute blood past the obstruction
    \item \textbf{Management of ascites:} diuretic therapy and a low-salt diet to control abdominal fluid
    \item \textbf{Management of varices:} endoscopic banding for bleeding oesophageal varices
    \item \textbf{Liver transplantation:} for people who develop advanced liver failure despite other measures
    \item \textbf{Specialist follow-up:} long-term care by a liver team with monitoring for disease progression
\end{itemize}

\section*{Complications}
The condition and its treatment can lead to serious complications.

\begin{itemize}
    \item Cirrhosis and progressive liver failure
    \item Bleeding from oesophageal varices
    \item Ascites and spontaneous bacterial peritonitis
    \item Bloodstream infection related to a shunt or stent
    \item Bleeding complications of anticoagulant therapy
    \item Chronic fatigue and malnutrition from impaired liver function
    \item Death from liver failure in severe, untreated cases
\end{itemize}

\section*{Prognosis and Outcomes}
Budd--Chiari syndrome is a serious condition, but the outlook has improved substantially with modern therapy. Many people respond well to anticoagulation combined with a stent or shunt and remain stable for many years. The prognosis is worse when diagnosis is delayed or when advanced liver failure or widespread hepatic fibrosis is already present; in these circumstances, liver transplantation offers good long-term survival. Outcomes depend on the underlying cause, the extent of hepatic involvement, and the speed of treatment.

\section*{Prevention and Self-Care}
The syndrome itself cannot always be prevented, but its risks can be modified.

\begin{itemize}
    \item If you have a known clotting disorder, follow specialist advice and attend regular appointments
    \item Discuss the risks of the combined oral contraceptive pill or hormone therapy with a doctor, especially if you have a thrombotic tendency
    \item Report new abdominal pain or swelling early, particularly if you take blood thinners
    \item Follow a low-salt diet and fluid advice if ascites develops
    \item Avoid alcohol, which damages the liver further
    \item Attend all liver clinic reviews and take prescribed medicines reliably
    \item Seek urgent help for bleeding, confusion, or rapid swelling of the abdomen
\end{itemize}

\section*{Sources and Further Reading}
The following organisations provide reliable information on Budd--Chiari syndrome and liver disease for patients and healthcare students.

\begin{itemize}
    \item NHS. \textit{Budd--Chiari syndrome: overview and treatment information.} https://www.nhs.uk/conditions/budd-chiari-syndrome/
    \item Mayo Clinic. \textit{Budd--Chiari syndrome: symptoms, causes, and treatment.} https://www.mayoclinic.org/
    \item MedlinePlus. \textit{Hepatic vein thrombosis (Budd--Chiari syndrome) information.} https://medlineplus.gov/
    \item MSD Manual. \textit{Budd--Chiari syndrome: professional overview.} https://www.msdmanuals.com/
    \item National Institute for Health and Care Excellence. \textit{Vascular and liver disease management guidance.} https://www.nice.org.uk/
    \item American Liver Foundation. \textit{Patient education on Budd--Chiari syndrome and liver health.} https://liverfoundation.org/
\end{itemize}
